\documentclass[10pt,a4paper]{amsart}
\usepackage[margin=19mm]{geometry}
\usepackage{amsmath,amssymb,mathtools}
\usepackage[scr=rsfs,cal=euler]{mathalfa}
\usepackage{xfrac}
\usepackage[unicode,hidelinks,bookmarks=false]{hyperref}
\newcommand{\ie}{i.e.\ }
\newcommand{\cf}{cf.\ }
\newcommand{\resp}{respectively\ }
\newcommand{\Subsection}{\subsection}
\providecommand{\nobreakdash}{\nobreak\hbox{-}}
\setlength{\emergencystretch}{2em}
\multlinegap0pt
\usepackage{bm}
\usepackage{bbm}
\usepackage{dsfont}
\usepackage{xspace}
\usepackage{cancel}
\usepackage{mathtools}
\usepackage{amsthm}
\makeatletter
\@ifundefined{thm}{\newtheorem{thm}{Theorem}}{}
\makeatother
\title[D-inverse constellations]{D-inverse constellations}
\author{Victoria Gould, Tim Stokes}
\date{Source: arXiv:2403.17282v2 (CC BY 4.0)}
\begin{document}
\maketitle

\noindent\emph{Source and licence.} D-inverse constellations. Version arXiv:2403.17282v2 (CC BY 4.0), distributed under the Creative Commons Attribution 4.0 International licence, as recorded in the arXiv licence metadata for this exact version and shown on the abstract page https://arxiv.org/abs/2403.17282v2. Full rights notice: licenses/arxiv-2403.17282-CC-BY-4.0.txt.\par\smallskip

\noindent\textit{Editorial scope.} Counted unit: the Thm whose source label is \texttt{invconsteg} in the article (source0.tex lines 558--560, source label \texttt{invconsteg}), delivered together with the article's own complete original proof of it (source0.tex lines 561--585, one \texttt{proof} environment of 25 lines). No mathematics is written, repaired or re-derived: statement and proof are the article's own text line for line. Required context retained uncounted: none. Every definition, display and label of those lines is retained unchanged. Omitted from this package and not redistributed: 1--557, 586--731. Nothing inside a retained excerpt is omitted or altered: every retained body is delivered exactly as the article's lines are written, so no omission receipt of any kind accompanies this package. Citations: the retained excerpts carry no citation command at all, so no bibliography entry of the article is transcribed and no reference is redistributed. Reference adaptation: none was needed and none was made; every reference inside the delivered unit resolves inside it, so no unresolved reference or citation is produced. Printed slips kept literal: no printed word, symbol, hypothesis or number of the counted statement or of its proof was altered. Layout and class: the article's own class is \texttt{amsart} with packages including comment, amsmath, latexsym, amssymb, eucal, enumitem, hyperref; the wrapper is the lane's pinned \texttt{amsart} editorial wrapper at 10pt on A4, which restates the theorem-like environments the retained text needs and adds the article's own macro definitions that the retained text needs (none), with extra wrapper packages (bm, bbm, dsfont, xspace, cancel, mathtools). The delivered numbering is the unit's own. Assets: no retained span contains a figure environment, a graphics-inclusion command, a PDF-page inclusion, a table or a TikZ picture; no image file is copied into this package, the package declares no asset dependency and no image file is redistributed with it. Licence: version arXiv:2403.17282v2 of this article is distributed under the Creative Commons Attribution 4.0 International licence, as recorded in the arXiv licence metadata for this exact version and shown on the abstract page https://arxiv.org/abs/2403.17282v2. A full rights notice is delivered at licenses/arxiv-2403.17282-CC-BY-4.0.txt. Characters: every retained excerpt is pure ASCII, so the delivered engine, XeLaTeX with its default Latin Modern text font, reports no missing character.
\begin{thm}  \label{invconsteg}
Suppose that $S$ is a semigroup containing non-empty $E\subseteq E(S)$, with $E\subseteq T\subseteq R_E(S)$. Then $\mathcal{I}_E(T)=(I_E(T),\cdot,D)$ is a D-inverse constellation with $(s',s)$ the D-inverse of $(s,s')$ for all $(s,s')\in I_E(S)$, if and only if $E$ is $T$-normal.
\end{thm}

\begin{proof}
Suppose first that $E$ is $T$-normal. Pick $(s,s')\in I_E(T)$.  Since $ss',s's\in E$, $(s,s')'=(s',s)\in I_E(T)$ also, and $(e,e)\in I_E(T)$ for all $e\in E$ since $e$ is an $E$-inverse of itself, so in particular, $D((s,s'))\in I_E(T)$. 

Suppose that $(s,s')\in I_E(T)$ is a right identity element. Certainly  
$(s',s)\in I_E(T)$ and the product $(s',s)\cdot (s,s')$ exists. It follows
that $(s',s)=(s's, s's)$, so that $s=s's\in E$.

Conversely, if $e\in E$ then $(e,e)\in I_E(T)$, and if $(x,x')\cdot (e,e)=(xe,ex')$ exists for some $(x,x')\in I_E(T)$, then $x=xee=xe$ and $x'=eex'=ex'$, so $(x,x')\cdot (e,e)=(xe,ex')=(x,x')$.  Hence $(e,e)$ is a right identity. It follows that the set of right identities of $\mathcal{I}_E(T)$ is given by
 \[\mathcal{RI}(I_E(T))=\{(e,e)\mid e\in E\}=\{D((s,s'))\mid (s,s')\in I_E(T)\}.\]

Choose $(s,s')\in I_E(T)$.  Then $ss'\in E\subseteq I_E(T)$, $(ss',ss')\in I_E(T)$, and  $D((s,s'))\cdot (s,s')$ exists (since $(ss')(ss')=ss'$ in $S$) and equals $(s,s')$.  If $(s,s')=(e,e)\cdot (s,s')=(es,s'e)$ for some right identity $(e,e)$, then $ess'=e$, $ss'e=e$, $es=s$ and $s'e=s'$. Hence $ss'=(es)(s'e)=(ess')e=ee=e$, showing uniqueness.

Now pick $(s,s'),(t,t')\in I_E(T)$ for which $(s,s')\cdot (t,t')$ exists,
so that $s=stt'$ and $tt's'=s'$.  Then $(s,s')\cdot (t,t')=(st,t's')$ and we must show that $(st,t's')\in I_E(T)$. We have that $$(st)(t's')(st)=(stt')(s'st)=s(s'st)=st$$ and $$(t's')(st)(t's')=t's'(stt')s'=t's'ss'=t's',$$ with $(st)(t's')=ss'\in E$.  But also, $(t's')(st)=t'(s's)t\in E$ by the $T$-normality of $E$.  Hence $(s,s')\cdot (t,t')=(ss',t't)\in  I_E(T)$, so this partial operation is well-defined, and moreover $D((s,s')\cdot (t,t'))=(stt's',stt's')=(ss',ss')=D((s,s'))$ from above.

If $(s,s'),(t,t'),(u,u')\in I_E(T)$ with $(s,s')\cdot (t,t')=(st,t's')$ and $(t,t')\cdot (u,u')=(tu,u't')$ both existing, then $s=stt', tt's'=s', tuu'=t$ and $ uu't'=t'$, so $s(tuu't')=stt'=s$ and $(tuu't')s'=tt's'=s'$, giving that $(s,s')\cdot ((t,t')\cdot (u,u'))$ exists.    

If $(s,s')\cdot((t,t')\cdot (u,u'))$ exists then it equals $(s,s')\cdot(tu,u't')=(stu,u't's')$ and we have $t=tuu'$, $uu't'=t'$, $s=s(tu)(u't')=stt'$, and $s'=(tu)(u't')s'=tt's'$, so $(s,s')\cdot (t,t')=(st,t's)$ exists. Further,  $(st)uu'=st$ and $uu't's'=t's'$, and so $((s,s')\cdot(t,t'))\cdot (u,u'))=(st,t's')\cdot (u,u')$ exists and is equal to $(stu,u't's')=(s,s')\cdot((t,t')\cdot (u,u'))$.

Hence, $(I_E(T),\cdot,D)$ is a constellation, and $D(I_E(T))=\{(e,e)\mid e\in E\}$.

 It is easy to check normality. Finally, for all $s\in I_E(T)$, we have $(s,s')\cdot (s',s)$ exists (since $ss's=s$ and $s'ss'=s'$ in $S$) and equals $(ss',ss')=D((s,s'))$, and similarly $(s',s)\cdot (s,s')$ exists and equals $D((s',s))$.  We have completed the verification that $\mathcal{I}_E(T)$ is a D-inverse constellation.

Conversely, suppose that $\mathcal{I}_E(T)$ is a D-inverse constellation as in the statement of the theorem. Let $e\in E$ and suppose that $s'$ is an $E$-inverse of $s$ for which $e=ess'=ss'e$.  It follows that $(e,e)\cdot (s,s')$ exists and equals $(es,s'e)$, so $s'e$ is an $E$-inverse of $es$.  In particular, $s'es=s'e(es)\in E$.  Hence $E$ is $T$-normal.  
\end{proof}

\end{document}
